%% kapa: exam=1 round=23 year=2012 period=1 form=A subject=economics subject_ko=경제원론 seq=14 number=54
%% source: Q-Net 감정평가사 기출문제 (cst00309) artlSeq=5052578 file=raw/1008/5052578/2012년 제23회 감정평가사 제1차시험 기출문제(1교시 A형).pdf
%% conversion: kordoc 4.24.51 -> tools/gen1.py (auto); review=proofread-sonnet; flags=none
\begin{question}{54}
동전을 던져 앞면이 나오면 9,000원을 따고 뒷면이 나오면 10,000원을 잃는 도박이 있다. 甲은 위험기피자, 乙은 위험애호자, 丙은 위험중립자인 경우 다음 설명으로 옳은 것은?
\begin{choices}
  \item 甲의 도박에의 참여 여부는 위험 기피도에 따라 결정될 것이다.
  \item 도박에 참여하는 대가로 500원을 준다 해도, 甲은 도박에 참여하지 않을 것이다.
  \item 丙은 이 도박에 반드시 참여할 것이다.
  \item 乙은 이 도박에 반드시 참여할 것이다.
  \item 앞면이 나올 때 따는 금액을 1,000원 올려 10,000원으로 하고, 뒷면이 나올 때 잃는 금액을 1,000원 내려 9,000원으로 하면 甲, 乙, 丙 모두 이 도박에 반드시 참여할 것이다.
\end{choices}
\end{question}
\answer{2}
